\chapter{Introduction}


Adversarial planning against non-cooperative agents has become a relevant area of research in real-world autonomous systems and robotics problems. Traditional optimal control frameworks to solve these kinds of problems tend to assume certain information strategies which can become jeopardized against an agent that's following a shifting, actively adversarial strategy coming from a versatile, intelligent opponent. 

In these cases, classical control architectures may prove to be brittle against these highly dynamic strategies. Inspired by this dilemma, we chose to develop a reinforcement-learning (RL) based approach for planning against these sorts of adversarial opponents. The benefit behind these RL-based architectures is that, when properly trained, these can produce a single, deterministic policy that can actively react and adapt to a wide arrange of scenarios it's been trained against. 

One of the arguments behind why research and industry have not widely adopted RL algorithms  within their autonomy/GNC work comes down to their  lack of formal guarantees for stability and convergence. This lack of guarantees can be attributed to the fact that these RL-based algorithms try to optimize against highly non-convex objectives while using gradients from noisy trajectories, approximate their value functions learned during trainings, are prone to overfit to the data distribution that they're trained on, and struggle to interpolate to data outside this distribution. In other words, these factors can lead to a noisy stochastic gradient ascent/descent process that finds locally optimal solutions instead of the globally optimal solution that may not be guaranteed to work in regimes outside the training dataset. In fact, the development of better RL algorithms that can account for these mishaps present in the current state-of-the-art RL is perhaps some of the most interesting research going on in the world of theoretical deep RL in this day and age. 

Regardless, algorithms such as Proximal Policy Optimization (PPO, \cite{schulman2017ppo}), an on-policy algorithm created by OpenAI, remains one of the most widely popular algorithm in RL research for engineering applications (\cite{francois2018introduction}). In simple terms, PPO is based on the principles of stochastic gradient ascent on a clipped surrogate objective function to be optimized. The benefits of these PPO-based architectures are that they're relatively simple to set up, they reuse each rollout batch on multiple gradient steps, are easy to parallelize, and, most importantly, provide much more stability than vanilla policy gradient methods. For controls based engineering applications, which is what we'll be studying in this work in this work, perhaps the biggest benefit of PPO is that it can be used for control continuous and discrete action spaces.



With all of the above in mind, we set off to study reinforcement learning (RL) framework to train, model, and rollout robust adversarial policies for orbital zero sum games in Low Earth Orbit (LEO). In order to better understand how these architectures could interact with classical autonomy methods, we were most interested in studying how this framework could learn while subject to a closed-loop aerospace guidance, navigation, and control architecture. For our own applications, we wanted to solve a few questions in this research:

\begin{itemize}
    \item Can an RL architecture use the PPO algorithm learn coherent policies for continuous control of orbital spacecraft in LEO in zero-sum adversarial scenarios?
    \item Can an RL architecture learn coherent policies under while interacting subject to the constrains of a classical controller in the loop?
    \item Can an RL architecture learn coherent policies if, instead of having access of the opponent's full state, it had partial observability and had to rely on measurements in-the-loop?
    \item Can an RL architecture interpolate outside of the training dataset even if it was never trained against some of that data?
    \item Can an RL architecture do all the above at once, provide convergent and high-performing policies despite the stochasticity?
\end{itemize}


With this in mind, we propose the Phased Adversarial Reinforcement Learning and Behavior Cloning (PARL+BC) method for training and rollout of RL-enabled Guidance Navigation and Control (GNC) aerospace applications. The method uses dense step per geometry-based rewards and phased opponent freezing to mitigate non-stationarity encountered in multi-agent training and aid towards learning stability.  PARL+BC learns while subject to the constrains of a classical velocity-saturating CBF-QP controller and, optionally, partial observability of the opponent's state via a bearings-only Extended Kalman Filter, which is achieved by using behavior cloning.



Our Monte Carlo evaluation results using the resulting policies show that PARL+BC produces highly performing policies in a constrained, zero-sum orbital asset defense game with equal spacecraft characteristics. The defender policies manage an aggregate success rate of up to 79.9\% across four learning phases in across different training scenarios of varying max actuation when having access to its opponents full state, where the opponent is also guided by an RL-trained policy. For the Ektended Kalman Filter-based policies, where neither agent has access to their opponents state and instead have to measure it via a bearings only EKF, the defender policies succeded in up to 74.7\% of the cases, a 6.51\% relative drop in performance relative to the fully observable policies. Furthermore, this was achieved while providing fast-rollout and compute inference times of 0.145 ms median per timestep on the full state policies and 3.91 ms per timestep on the EKF-based policies.

Our findings open the way for the development of highly compute efficient RL-based method for autonomous systems research and application, and a discussion of potential directions for future work is provided.

For guidance on the terminology used throughout this writing, please refer to appendix \ref{appendix:glossary} which serves as a glossary of some of the key terms used throughout this research.